\documentclass[openany,a5paper]{utbook}

%% heading
\usepackage{ctexheading}
\usepackage{zhnumber}
\ctexset{
  chapter = {
    name        = 古文觀止卷之 ,
    number      = \zhnum{chapter} ,
    format      = \huge\bfseries ,
    nameformat  = {} ,
    titleformat = {} ,
    aftername   = \quad ,
    beforeskip  = 0 pt ,
    afterskip   = 20 pt ,
    afterindent = true ,
    pagestyle   = main ,
  } ,
  section = {
    numbering   = false ,
    indent      = 42 pt ,
    afterindent = true ,
  } ,
}
\zhnumsetup{style=Traditional}

%% margin
\usepackage{titleps}
\makeatletter
    \newpagestyle{main}{%
      \widenhead{30pt}{30pt}%
      % TODO -- Set the headings of even pages
      \sethead[][][%
      {\if@mainmatter
        \raisebox{\dimexpr-\height-\headsep\relax}[0pt]{%
          \hbox to \textheight{%
            \tate\hspace{3zh}\CTEXthechapter\hfill
            \zhdigits{\arabic{page}}\hspace{4zh}}}%
        \fi}]{%
        \if@mainmatter
        \raisebox{\dimexpr-\height-\headsep\relax}[0pt]{%
          \hbox to \textheight{%
            \tate\hspace{3zh}\CTEXthechapter\hfill
            \zhdigits{\arabic{page}}\hspace{4zh}}}%
        \fi}{}{}
    }\makeatother
\makeatletter
\pagestyle{main}

%% proper name & book name marks
\usepackage[normalem]{ulem}
\newcommand\ProperName[1]{\hskip.1zw\uline{\kern-.1zw#1\kern-.1zw}\kern.1zw}
\newcommand\PlaceName[1]{\edef\tempdepth{\the\ULdepth}\setlength{\ULdepth}{.75zw}\hskip.1zw\uuline{\kern-.1zw#1\kern-.1zw}\kern.1zw\setlength{\ULdepth}{\tempdepth}}
\newcommand\smallwave{\bgroup \markoverwith{\lower3.5pt\hbox{\font\tempfont=lasy5 at .2zw\tempfont \char58}}\ULon}
\newcommand\BookTitle[1]{\hskip.1zw\smallwave{\kern-.1zw#1\kern-.1zw}\kern.1zw}

%% fonts
\usepackage[dvipdfmx,bookmarksnumbered]{hyperref}
\AtBeginShipoutFirst{\special{pdf:tounicode UTF8-UTF16}}

\DeclareKanjiFamily{JY2}{song}{}
\DeclareFontShape{JY2}{song}{m}{n}{<->upschrm-h}{}
\DeclareFontShape{JY2}{song}{bx}{n}{<->ssub*hei/m/n}{}
\DeclareKanjiFamily{JY2}{hei}{}
\DeclareFontShape{JY2}{hei}{m}{n}{<->upschgt-h}{}
\DeclareFontShape{JY2}{hei}{bx}{n}{<->ssub*hei/m/n}{}
\DeclareKanjiFamily{JT2}{song}{}
\DeclareFontShape{JT2}{song}{m}{n}{<->upschrm-v}{}
\DeclareFontShape{JT2}{song}{bx}{n}{<->ssub*hei/m/n}{}
\DeclareKanjiFamily{JT2}{hei}{}
\DeclareFontShape{JT2}{hei}{m}{n}{<->upschgt-v}{}
\DeclareFontShape{JT2}{hei}{bx}{n}{<->ssub*hei/m/n}{}
\AtBeginShipoutFirst{%
  \special{pdf:mapline upstsl-h UniGB-UTF16-H simsun.ttc}%
  \special{pdf:mapline upstsl-v UniGB-UTF16-V simsun.ttc}%
  \special{pdf:mapline upstht-h UniGB-UTF16-H simhei.ttf}%
  \special{pdf:mapline upstht-v UniGB-UTF16-V simhei.ttf}}

\renewcommand\mcdefault{song}
\renewcommand\gtdefault{hei}

%% misc layout settings
\renewcommand{\thefootnote}{〔\zhdigits{\arabic{footnote}}〕}
\setlength\parindent{2zw}

\begin{document}
\mainmatter
\chapter{}
\section[鄭伯克段于鄢]{\ProperName{鄭伯}克\ProperName{段}于\PlaceName{鄢}}

初，\ProperName{鄭武公}娶于\PlaceName{申}\footnote{注一}，曰\ProperName{武姜}，生\ProperName{莊公}及\ProperName{共叔段}。\ProperName{莊公}寤生，驚\ProperName{姜氏}，故名曰\ProperName{寤生}，遂惡之。愛\ProperName{共叔段}，欲立之，亟請於\ProperName{武公}，\ProperName{公}弗許。

及\ProperName{莊公}即位，為之請\PlaceName{制}。\ProperName{公}曰：「\PlaceName{制}，巖邑也。\ProperName{虢叔}死焉，佗邑唯命\footnote{再注}。」請京，使居之，謂之\ProperName{京城大叔}。

\ProperName{祭仲}曰：「都城過百雉，國之害也。先王之制：大都，不過參國之一；中，五之一；小，九之一。今京不度，非制也，君將不堪。」\ProperName{公}曰：「\ProperName{姜氏}欲之，焉辟害？」對曰：「\ProperName{姜氏}何厭之有？不如早為之所，無使滋蔓！蔓，難圖也。蔓草猶不可除，況君之寵弟乎？」\ProperName{公}曰：「多行不義，必自斃，子姑待之！」

既而\ProperName{大叔}命\PlaceName{西鄙}、\PlaceName{北鄙}貳於己。\ProperName{公子呂}曰：「國不堪貳，君將若之何？欲與\ProperName{大叔}，臣請事之；若弗與，則請除之，無生民心。」\ProperName{公}曰：「無庸，將自及。」

大叔又收貳以為己邑，至於廩延。子封曰：「可矣！厚將得眾。」公曰：「不義不暱，厚將崩。」

大叔完聚，繕甲兵，具卒乘，將襲鄭，夫人將啟之。公聞其期曰：「可矣。」命子封帥車二百乘以伐京，京叛大叔段。段入于鄢，公伐諸鄢。五月辛丑，大叔出奔共。

書曰：「鄭伯克段于鄢。」段不弟，故不言弟。如二君，故曰克。稱鄭伯，譏失教之。謂之鄭志。不言出奔，難之也。

遂寘姜氏于城潁，而誓之曰：「不及黃泉，無相見也。」既而悔之。潁考叔為潁谷封人，聞之。有獻於公，公賜之食。食舍肉，公問之。對曰：「小人有母，皆嘗小人之食矣。未嘗君之羹，請以遺之。」公曰：「爾有母遺，繄我獨無！」潁考叔曰：「敢問何謂也？」公語之故，且告之悔。對曰：「君何患焉。若闕地及泉，隧而相見，其誰曰不然？」公從之。

公入而賦：「大隧之中，其樂也融融。」姜出而賦：「大隧之外，其樂也泄泄！」遂為母子如初。

君子曰：「潁考叔，純孝也，愛其母，施及莊公。\BookTitle{詩}曰：『孝子不匱，永錫爾類』{，}其是之謂乎？」


\section{周鄭交質}
鄭武公、莊公為平王卿士。王貳于虢，鄭伯怨王。王曰：「無之。」故周鄭交質：王子狐為質於鄭，鄭公子忽為質於周。

王崩，周人將畀虢公政。四月，鄭祭足帥師取溫之麥。秋，又取成周之禾。周鄭交惡。

君子曰：「信不由中，質無益也。明恕而行，要之以禮，雖無有質，誰能間之？苟有明信：澗溪沼沚之毛，蘋蘩薀藻之菜，筐筥錡釜之器，潢汙行潦之水，可薦於鬼神，可羞於王公。而況君子結二國之信，行之以禮，又焉用質？\BookTitle{風}有\BookTitle{采}\BookTitle{蘩}、\BookTitle{采蘋}，\BookTitle{雅}有\BookTitle{行葦}、\BookTitle{泂酌}，昭忠信也。」


\section{石碏諫寵州吁}
衞莊公娶于齊東宮得臣之妹，曰莊姜，美而無子。衞人所為賦\BookTitle{碩人}也。又娶于陳，曰厲媯，生孝伯，蚤死。其娣戴媯，生桓公，莊姜以為己子。

公子州吁，嬖人之子也，有寵而好兵，公弗禁。莊姜惡之。

石碏諫曰：「臣聞愛子，教之以義方，弗納於邪。驕、奢、淫、佚，所自邪也。四者之來，寵祿過也。將立州吁，乃定之矣；若猶未也，階之為禍。夫寵而不驕，驕而能降，降而不憾，憾而能眕者，鮮矣。且夫賤妨貴，少陵長，遠間親，新間舊，小加大，淫破義，所謂六逆也；君義，臣行，父慈，子孝，兄愛，弟敬，所謂六順也。去順效逆，所以速禍也。君人者，將禍是務去；而速之，無乃不可乎？」弗聽。

其子厚與州吁遊，禁之，不可。桓公立，乃老。


\chapter{}
\section{鄭子家告趙宣子}
晉侯合諸侯于扈，平宋也。於是晉侯不見鄭伯，以為貳于楚也。 鄭子家使執訊而與之書，以告趙宣子，曰：「寡君即位三年，召蔡侯而與之事君。九月，蔡侯入于敝邑以行；敝邑以侯宣多之難，寡君是以不得與蔡侯偕。十一月，克減侯宣多，而隨蔡侯以朝於執事。

十二年六月，歸生佐寡君之嫡夷，以請陳侯於楚，而朝諸君。十四年七月，寡君又朝以蕆陳事。十五年五月，陳侯自敝邑往朝於君。往年正月，燭之武往朝夷也。八月，寡君又往朝。以陳、蔡之密邇於楚，而不敢貳焉，則敝邑之故也。雖敝邑之事君，何以不免？在位之中，一朝于襄，而再見於君。夷與孤之二三臣，相及於絳。雖我小國，則蔑以過之矣。

今大國曰：『爾未逞吾志。』敝邑有亡，無以加焉！古人有言曰：『畏首畏尾，身其餘幾？』又曰：『鹿死不擇音。』小國之事大國也，德，則其人也；不德，則其鹿也。鋌而走險，急何能擇？命之罔極，亦知亡矣。將悉敝賦，以待於鯈，唯執事命之。

文公二年，朝於齊；四年，為齊侵蔡，亦獲成於楚。居大國之閒，而從於強令，豈其罪也？大國若弗圖，無所逃命。」

晉鞏朔行成於鄭，趙穿、公婿池為質焉。

\end{document}
